%% SYNTHETIC TEST UNIT generated by writing/tools/make_outline_book.py from FINAL_OUTLINE_v2 -- not book text
\chapter{喷注、事件形状、非全局对数与超领头对数}\label{ch:jets}
本单元为测试用的合成文本（由 FINAL\_OUTLINE\_v2 自动生成），陈述内容从略。

\section{第 1 节（测试）}\label{sec:jets-irsafe}

\begin{definition}[红外安全的喷注定义与事件形状]\label{def:jets-algorithms}
（测试用合成陈述：内容从略。）
\uses{def:ir-safe}
\end{definition}

\begin{proposition}[Sterman–Weinberg 两喷注截面 \tagR]\label{prop:jets-sw}
（$R_{\mathrm{PT}}$：在重整化微扰论中对红外正规化的部分子态严格。）
（测试用合成陈述：内容从略。）
\uses{def:jets-algorithms, prop:ee-oneloop, prop:qcd-soft-collinear-limits}
\end{proposition}

\begin{example}[$O(\alpha_s)$ 的 Sterman–Weinberg 截面]\label{ex:jets-sw}
（测试用合成陈述：内容从略。）
\uses{prop:jets-sw}
\end{example}

\section{第 2 节（测试）}\label{sec:jets-hadronization}

\begin{corollary}[红外安全喷注截面的有限性 \tagA]\label{cor:jets-ir-finite}
（测试用合成陈述：内容从略。）
\uses{thm:ir-safe-finite, def:jets-algorithms}
\end{corollary}
\begin{proofsketch}测试。\end{proofsketch}
\paragraph{地位与缺口}测试：见文献。

\begin{proposition}[领头阶推力分布 \tagR]\label{prop:jets-thrust-lo}
（$R_{\mathrm{PT}}$：在重整化微扰论中对红外正规化的部分子态严格。）
（测试用合成陈述：内容从略。）
\uses{def:jets-algorithms}
\end{proposition}

\begin{conjecture}[强子化修正的普适性 \tagC]\label{conj:jets-hadronization}
（测试用合成陈述：内容从略。）
\uses{thm:scet-thrust}
\end{conjecture}
\paragraph{地位与缺口}测试：见文献。

\section{第 3 节（测试）}\label{sec:jets-ngl}

\begin{proposition}[非全局对数的存在 \tagR]\label{prop:jets-ngl}
（$R_{\mathrm{PT}}$：在重整化微扰论中对红外正规化的部分子态严格。）
（测试用合成陈述：内容从略。）
\uses{def:jets-algorithms}
\end{proposition}

\begin{theorem}[BMS 方程 \tagA]\label{thm:jets-bms}
（测试用合成陈述：内容从略。）
\uses{prop:jets-ngl}
\end{theorem}
\begin{proofsketch}测试。\end{proofsketch}
\paragraph{地位与缺口}测试：见文献。

\begin{proposition}[超出领头对数与有限 $N_c$ \tagP]\label{prop:jets-ngl-beyond}
（测试用合成陈述：内容从略。）
\uses{thm:jets-bms}
\end{proposition}
\paragraph{地位与缺口}测试：本陈述的地位由\cref{prop:front-ngl}决定。

\section{第 4 节（测试）}\label{sec:jets-sll}

\begin{theorem}[喷注过程的有效场论因子化 \tagA]\label{thm:jets-ngl-factorization}
（测试用合成陈述：内容从略。）
非微扰地位：\cref{asm:intro-bridge}。
\uses{thm:scet-thrust, prop:jets-ngl, asm:intro-bridge}
\end{theorem}
\begin{proofsketch}测试。\end{proofsketch}
\paragraph{地位与缺口}测试：本陈述的地位由\cref{thm:scet-thrust}决定。

\begin{proposition}[超领头对数 \tagP]\label{prop:jets-sll}
（测试用合成陈述：内容从略。）
\uses{prop:sudakov-coulomb, thm:sudakov-ir-poles}
\end{proposition}
\paragraph{地位与缺口}测试：本陈述的地位由\cref{prop:break-spacelike}决定。
